\documentclass[11pt,a4paper]{article}
\usepackage[T1]{fontenc}
\usepackage[utf8]{inputenc}
\usepackage[ngerman]{babel}
\usepackage[margin=2.5cm]{geometry}
\usepackage[hidelinks]{hyperref}
\usepackage{array}
\setlength{\parindent}{0pt}
\setlength{\emergencystretch}{2em}
\renewcommand{\arraystretch}{1.25}
\sloppy
\begin{document}
\shorthandoff{"}
\begin{center}
  {\Large\textbf{Webtechnologien -- Übungsblatt 6}}
\end{center}
\section*{Aufgabe 3: Database-Backed Product List API}
\textbf{Aufgabe:} Produkte aus PostgreSQL laden und die Beschreibung
über PUT ändern. Fehlende Produkte, ungültige Eingaben und
Datenbankfehler sollen passende HTTP-Antworten liefern.

\medskip
\textbf{Ergebnis:} Ergänzt sind \texttt{server.js}, \texttt{pool.js} und
\texttt{config.json}. Die Vorlagen bleiben als \texttt{server\_ORIGINAL.js},
\texttt{pool\_ORIGINAL.js} und \texttt{config\_ORIGINAL.json} erhalten.
\begin{itemize}
  \item \texttt{config.json}: Der überflüssige Beistrich nach dem
    Datenbankobjekt ist entfernt, damit gültiges JSON entsteht.
    Host: \texttt{localhost}; Benutzer und Passwort: \texttt{postgres};
    Datenbank: \texttt{webtechProducts} (lokale Testumgebung).
  \item \texttt{pool.js}: \texttt{new Pool} verwaltet wiederverwendbare
    Verbindungen. Die Werte kommen aus der Konfiguration.
    Deren Schlüssel heißt \texttt{db}, node-postgres erwartet dafür
    \texttt{database}. \texttt{pool.on('error', ...)} meldet einen
    Verbindungsabbruch, statt den Server abstürzen zu lassen.
    \texttt{module.exports} stellt den Pool bereit.
  \item GET \texttt{/products}: \texttt{SELECT id FROM products} lädt
    die IDs; \texttt{result.rows.map} bildet die Antwortliste.
    Ohne \texttt{ORDER BY} ist die Reihenfolge nicht garantiert.
  \item GET \texttt{/product/:id}: Die Abfrage lädt \texttt{id},
    \texttt{title}, \texttt{description} und \texttt{price}.
    \texttt{\$1} ist ein Platzhalter; die ID wird getrennt als Wert
    übergeben und kann dadurch keinen SQL-Befehl einschleusen.
    Eine leere Trefferliste bedeutet 404.
  \item PUT prüft zuerst den Typ, einen nicht leeren Inhalt und höchstens
    100 Zeichen. Nur Leerzeichen sind ebenfalls ungültig.
    \texttt{UPDATE ... RETURNING} ändert die Beschreibung und gibt die
    ganze neue Zeile zurück. \texttt{rowCount === 0} bedeutet 404.
    Auch hier werden Beschreibung und ID getrennt übergeben.
  \item \texttt{bodyParser.json()} macht den JSON-Text als
    \texttt{req.body} zugänglich. \texttt{async/await} wartet auf die
    Abfrage; jede Route fängt Fehler mit \texttt{try/catch} ab,
    protokolliert sie mit \texttt{console.error} und sendet 500.
\end{itemize}

\textbf{Starten:} PostgreSQL starten. Im Ordner \texttt{aufgabe3} zuerst einmal
\texttt{npm install} ausführen. Import im Unterordner:
\begin{quote}
\texttt{cd db\_import}\\
\texttt{psql -U postgres -f createDB.sql}\\
\texttt{cd ..}
\end{quote}
Die bereits korrigierte \texttt{createDB.sql} verwendet
\texttt{CREATE DATABASE "webtechProducts"} und
\texttt{ALTER DATABASE "webtechProducts"}. Ohne Anführungszeichen macht
PostgreSQL den Namen zu \texttt{webtechproducts}; das anschließende
\texttt{\textbackslash connect webtechProducts} scheitert dann und die
Tabelle wird nicht in der vorgesehenen Datenbank angelegt.
Die unveränderte Vorlage bleibt unter
\texttt{db\_import/createDB\_ORIGINAL.sql} erhalten.
Falls Windows die Locale \texttt{en\_US.UTF-8} ablehnt, beim dortigen
Import die Teile \texttt{LC\_COLLATE} und \texttt{LC\_CTYPE} entfernen.
In pgAdmin die Datenbank und unter \texttt{public} die Tabelle
\texttt{products} mit vier Zeilen kontrollieren.

Mit \texttt{node server.js} starten und
\url{http://localhost:3000/products} öffnen (Port 3000 muss frei sein).
PUT in Postman/RESTer an
\url{http://localhost:3000/product/sony24105} senden:
Header \texttt{Content-Type: application/json}, Body
\texttt{\{"description": "Neue Beschreibung"\}}.
Danach GET auf dieselbe URL ausführen oder die Zeile in pgAdmin prüfen.
Strg+C beendet den Server.

\medskip
\textbf{Test:} Mit Node.js v24 und PostgreSQL 17 (frischer Import) gestartet, Anfragen mit \texttt{curl}; alle Fälle liefern genau diese Antworten. Den Datenbankfehler
habe ich durch Stoppen von PostgreSQL ausgelöst.
\begin{center}
\begin{tabular}{|>{\raggedright\arraybackslash}p{5.2cm}|>{\raggedright\arraybackslash}p{8.5cm}|}
\hline
\textbf{Fall} & \textbf{Ergebnis} \\ \hline
GET \texttt{/products} & 200; \texttt{\{"productIds":["sonyAIV", "sony24105", "NiconDX16400", "tamron1724"]\}}; Reihenfolge darf abweichen \\ \hline
GET \texttt{/product/sony24105} & 200; \texttt{\{"id":"sony24105", "title":"Sony 24-105mm lens", "description":"Best Sony travel zoom", "price":1300\}} \\ \hline
GET \texttt{/product/nope} & 404; \texttt{\{"error":"product not found"\}} \\ \hline
PUT \texttt{/product/sony24105}, Body wie oben & 200; \texttt{\{"id":"sony24105", "title":"Sony 24-105mm lens", "description":"Neue Beschreibung", "price":1300\}} \\ \hline
PUT \texttt{/product/nope}, gültige Beschreibung & 404; \texttt{\{"error":"product not found"\}} \\ \hline
PUT: Beschreibung fehlt, kein String, leer, nur Leerzeichen oder mehr als 100 Zeichen & 400; \texttt{\{"error":"description must be a non-empty string with at most 100 characters"\}} \\ \hline
Datenbankfehler bei GET \texttt{/products}, GET \texttt{/product/:id} oder gültigem PUT & 500; \texttt{\{"error":"database error"\}}; Ursache im Serverterminal \\ \hline
\end{tabular}
\end{center}

\medskip
{\small \textbf{Hilfsmittel:} Claude Code (Claude Opus 5.5) und OpenAI Codex (GPT-6.1-Sol) für Erstellung, Prüfung und Dokumentation.}
\end{document}
